\subsection{Область допустимых значений крутящего момента на выходном валу передаточного механизма}

Согласно требованиям технического задания максимальный крутящий момент на выходном валу передаточного механизма не должен превышать
\[
|M_m(t)| \le 150 \text{ Н}\cdot\text{м}.
\]

Для определения момента на выходном валу воспользуемся выражением, полученным при составлении математической модели механической части:
\begin{equation}
    M_{\text{м}}(t)
    =
    b
    \left(
        \dot{\varphi}_{\text{рпр}}(t)
        -
        \dot{\varphi}_{\text{м}}(t)
    \right)
    +
    c
    \left(
        \varphi_{\text{рпр}}(t)
        -
        \varphi_{\text{м}}(t)
    \right).
\end{equation}

Для каждой точки области устойчивости вычислялось максимальное значение момента
\[
|M_m|_{\max}=\max_t |M_m(t)|.
\]

Далее среди устойчивых режимов были выделены значения коэффициентов регулятора \(k_{\text{п}}\) и \(k_{\text{и}}\), удовлетворяющие ограничению
\[
|M_m|_{\max} \le 150 \text{ Н}\cdot\text{м}.
\]


На рисунке~\ref{fig:pdf/moment_region} красным цветом показана область устойчивых значений коэффициентов регулятора, а зелёным цветом — часть области устойчивости, удовлетворяющая ограничению по максимальному крутящему моменту.

\screenshot[0.8]{pdf/moment_region}{Область допустимых коэффициентов регулятора по условию ограничения максимального момента}

Из полученного графика видно, что при увеличении коэффициента интегральной составляющей \(k_{\text{и}}\) максимальное значение момента возрастает быстрее, чем при изменении коэффициента \(k_{\text{п}}\). Вследствие этого часть устойчивых режимов не удовлетворяет ограничению
\[
|M_m|_{\max} \le 150 \text{ Н}\cdot\text{м}.
\]

Таким образом, допустимая область параметров регулятора определяется не только условием устойчивости системы, но и ограничением на максимальный крутящий момент на выходном валу передаточного механизма.
